\chapter{Sūtrasthāna 6: Seasonal living}

% K 006r IMG_0141.JPG" line 1.
% H dscn2986 fol 010.jpg upper line 6
% N A 1267_11_06.jpg upper line 2

\section{Translation}


\begin{translation}    
  
\item [1] 
Now, we will explain seasonal living (\emph{ṛtucaryā}).
% K has ṛtucarya, which is also possible as a neuter noun (cf. bhikṣācarya, brahmacarya, dauścarya, etc.) But the other witnesses and vulgate have rṭucaryā. 


\item [2] 
Time is indeed venerable, % kālo hi bhagavān
% Ḍalhaṇa: hi upadarśane, 
self-existent, and % svayaṃbhur
without beginning, middle or end. % anādimadhyanidhano
%
The loss or abundance of the tastes % rasasya madhurādeḥ, rasaśabdenātra rasavaddravyam ucyate'bhedopacārād ity anye
and people's life and death % jīvitamaraṇe ca
depend on it. % atra ... āyatte (agreeing with jīvitamaraṇe)
It is called time (\emph{kāla}) because % iti kālaḥ
it does not disappear % sa ... na līyate
even for a brief moment (\emph{kalā}). % sūkṣmām api kalān: accusative of extent?: Ḍalhaṇa saḥ kālaḥ, sūkṣmām api stokām api, kalāṃ bhāgaṃ, na līyate gatimattvāt śliṣṭo na bhavati; 
It is called time because % % iti kālaḥ
it gathers or impels living beings.%
    \footnote{According to Ḍalhaṇa, this derivation (\emph{nirukti}) of the word time (\emph{kāla}) is based on its action of making all beings into one mass by destroying them (\emph{saṃharaṇād ekarāśīkaroti}) or impelling them towards death (\emph{kālayati mṛtyusamīpaṃ nayati}).}

\item [3] 
Dependent on time as the year, % tasya samvatsarātmano : Cf. Macdonald Sanskrit Grammar p. 193: the genitive in the sense of " dependent on, belongoing or attached to."
the venerable sun % bhagavān ādityo 
produces  % karoti
through the particularity of its motion % gativiśeṣeṇa
the divisions of time % pravibhāgaṃ
into a \emph{nimeṣa}, %twinkling, % nimeṣa,
\emph{kāṣṭhā}, % jiffy, % kāṣṭhā,
\emph{kalā}, %short interval, %kalā,
forty-eight minute interval (\emph{muhūrta}), %muhūrta,
a day and night, %ahorātra
fortnight, % pakṣa
month, % māsa 
season, %rtv 
solar course (half a year), %ayana
a year, %samvatsara
and \emph{yuga} (five years). % yuga.

\item [4] 
Among those, % tatra
a \emph{nimeṣa} is the duration of a quick blink of the eye. %
Fifteen \emph{nimeṣa}s is a \emph{kāṣṭhā}. %
Thirty \emph{kāṣṭhā}s is a \emph{kalā}. %
Twenty and one-tenth \emph{kalā}s make a forty-eight minute interval. %
Thirty forty-eight minute intervals is a day and night. %
Fifteen days and nights is a fortnight. %
There are two fortnights: dark and bright. Together, the two make a month. 

% a twinkling is the duration of a quick blink of the eye. % nimeṣa
% Fifteen twinklings is a jiffy. % 15 nimeṣa = kāṣṭhā
% Thirty jiffies is an interval. % 30 kāṣṭhās = kalā
% Twenty and a tenth intervals is forty-eight minutes. % 20 and a tenth
% Thirty forty-eight minute periods is a day and night.

\item [5] 
In regard to that, % tatra
the twelve months beginning with Māgha make a year. %
Considering each season is two months long, there are six seasons. %
And these are late winter, %śiśira 
spring, % vasanta
summer, % grīṣma 
the monsoon season, % varṣā 
autumn, % śarad 
and early winter. % hemanta
%
Of those, the dewy season is the months Tapas and Tapasya. Spring is Madhu and Mādhava. Summer is Śuci and Śukra. The rainy season is Nabho and Nabhasya. Autumn is Iṣa and Ūrja, and winter is Sahaḥ and Sahasya.
    % should we adopt conventional sandhi in the edition  fo sahaḥsahasyau or stick with K and N (H has sahaḥsahasyau)?
    % should iti be translated here?

\item [6] 
These seasons are characterised by cold, heat, rain and wind. %
Since the sun and moon are the agents of dividing time, % candrādityayoḥ kālavibhāgakaratvād
there are two solar courses: % ayana
the southern and northern. % I assume that this refers to dakṣiṇāyana — the southern course of the Sun and uttarāyaṇa — the northern course of the Sun.
Of the two, the southern is the rainy season, autumn, and winter. During those seasons, the venerable moon is strong, % āpyāyate: to swell, increase; to grow larger or fat or comfortable; to thrive; to become full or strong; to abound. Obviously the moon is waxing and waning during these seasons, so perhaps āpyāyate implies the moon's influence or power is strong or increasing?
and the dominant tastes are sour, salt, and sweet. % amlalavaṇamadhurāś
The strength of all beings also grows continuously. % uttarottarañ : exponentially (maybe too strong?)
The northern course is late winter, spring, and summer. % śiśiravasantagrīṣmāḥ
During these seasons, the venerable sun is strong. %
The strongest tastes are sharp, spicy, and astringent. % Check against previous translations
The strength of beings also declines continuously.  

\item [7] 
And on this there is the following verse:
\begin{quote}
The moon moistens the earth, then the sun dries it.  Depending on both, % saṃśritya (having resorted to) is a bit difficult to construe here 
the wind sustains living creatures. 
\end{quote}

\item [8] 
Now, as it well known, both solar courses % ayane (cf. the vulgate ayane dve) 
together % yugapat
are the year. % saṃvatsaro bhavati
These two courses %  te dve ayane
constitute % supplied
each of the following five years: %  ity evaṃ pañca pañca varṣāṇi. not sure how to render pañca pañca... a set of? 
Varṣa, Saṃvatsara, Parivatsara, Iḍāvatsara, and Vatsara. %
The five are designated % saṃjñāṃ labhante: acquire the designation
a \emph{yuga}. %
This time period, beginning with a \emph{nimeṣa} and ending with a \emph{yuga}, is revolving like a wheel, and is thus called, `the wheel of time (\emph{kālacakra}).' 

\item [9] 
In this way, night during the southern solar course is explained. % not sure what this means, and it has been removed from the vulgate.

\item [10] 
In this tradition, % Ḍalhaṇa:  iha asminn adhyāye, but he thinks the issue here is why the seasons are being mentioned again in this chapter (asminnadhyāye punaḥ kimarthaṃ varṣādayaḥ ṛtava ityāha- ) and his answer is doṣopacayādyartham
there are six seasons, which are the monsoon, autumn, early winter, spring, summer, and late winter, because of the accumulation, aggravation and pacification of the humours. 
    % is pacification the best translation of praśama?
They are thus explained % te ... evaṃ vyākhyātāḥ
by two-month periods, % dvimāsikena
beginning with Bhādrapada, % Is there a translation for Bhādrapada? The dictionary says, "N. of the 3rd and 4th lunar asterisms."
% not the difference in syntax with the vulgate (bhādrapadādyair dvimāsikena vs bhādrapadādyena dvimāsikena)
as follows: % tad yathā ... iti
Bhādrapada and Āśvayuja constitute the rainy season; Kārttika and Mārgaśīrṣa, autumn; Pauṣa and Māgha, early winter; Phālguna and Caitra, spring; Vaiśākha and Jyeṣṭha, summer; and Āṣāḍha and Śrāvaṇa, late dewy season.
    % I've kept the daṇḍas in the edition to avoid changing sandhi. 

\item [11]
Here, in the rainy seasons, medicines are fresh, have low potency, and, when taken as food, % āhāratvam upagatā (lit: having become the state of food)
produce sour digestion.\footnote{%
    Ḍalhaṇa \citep[26]{vulgate} interprets ``they are burnt'' (\emph{vidahyante}) as ``they produce sour digestion (\emph{amlapākam upayānti}).} % Ḍalhaṇa:  vidahyante amlapākam upayānti
Water is also stirred up % apraśāntāḥ
and contains dirt and impurities. % kṣitimalaprāyās
However, % tu
When they are consumed % tāḥ ... upayujyamānāḥ
while the sky is covered with clouds % nabhasi meghāvatate
and the earth has been moistened with water, % jalapraklinnāyāṃ bhūmau: is soaked too strong?
they produce sour digestion % tāḥ vidahyante
for people whose bodies are moist % klinnadehānāṃ dehināṃ
and whose digestive fire has been obstructed by the cold, wind, and rain. % śīta-vāta-varṣa-viṣṭambhitāgnīnām
Owing to the sour digestion, they cause an accumulation of bile.%
    % All Nepalese witnesses have savidāhāt, but I can't see how that works here. So I have reverted to the reading of the vulgate. 


In autumn, %śaradi
when the sky has very few clouds % praviralameghe viyati
and the mud is drying up, % upaśuṣyati (PAP) paṅke
accumulation [of bile] is dissolved by the rays of the sun % sa sañcayaḥ ... arkakiraṇapravilāyitaḥ
and produces bile-related diseases. %
Owing to the change of seasons, %kālapariṇāmāt
those medicines whose potency has matured % pariṇatavīryāḥ
become powerful. 
Also, water is still, viscous % snigdhā
and excessively heavy. % atyarthaṃ gurvyaḥ
Being consumed, they produce sweet digestion % avidagdhāḥ ? 
for people whose bodies are stiffened by the cold % satuṣāraupaṣṭambhitadehānāṃ
because of the weak rays of the sun.% mandakiraṇatvād bhānos 
    \footnote{Ḍalhaṇa \citep[26]{vulgate} interprets ``they are not burnt'' (\emph{avidagdhāḥ}) as ``they produce sweet digestion (\emph{madhurapākam upagatāḥ}). }
Owing to their viscosity, heaviness and filmy quality % upalepitvāc ca ?
they create an accumulation of phlegm. %

In Spring, the accumulation [of phlegm] is dissolved by the rays of the sun and produces phlegm-related diseases.  %śleṣmikān vyādhīn

In Summer, those medicines % tā evauṣadhyaḥ
become insubstantial, dry, and extremely light,\footnote{%
    Ḍalhaṇa \citep[26]{vulgate} understands ``insubstantial'' (niḥsāra) in terms of their weight (\emph{sāro ’tra gauravam}), which seems to emphasise Śuśruta's further description of these medicines as ``extremely light'' (\emph{atimātralaghu}).} 
and [so too] water.\footnote{%
    Ḍalhaṇa \citep[26]{vulgate} states that, in summer, water possess the same qualities as medicine (\emph{na kevalam oṣadhaya evaṃvidhā bhavanti āpaś ca bhavantīty arthaḥ}). In this passage the  the water's qualities are not stated separately.}
When they are consumed, % upayujyamānās
they produce an accumulation of wind because of their dryness and lightness, % raukṣyāl lāghavāc ca
for people whose bodies have been dried out by the heat of the sun.% sūryapratāpopaśoṣitadehānān dehināṃ
    \footnote{% 
    In K, the following passage has been inserted midway through this sentence: ``Here, in autumn and winter, day and night are the same, and two intercalary months occur simultaneously. During the northern solar course, the length of the day increases by twenty \emph{kāṣṭhā}s and half a \emph{nimeṣa}.'' This statement breaks the syntax of the surrounding sentence and does not fit the context of this section, which concerns the effects of the seasons on medicines and water. It is likely, therefore, that it was originally composed as a marginal note that was mistakenly incorporated into the main text at some stage in the transmission. In K, this insertion has been marked with round brackets by a later hand and does not appear in N or H.} %
  
In the monsoon, % prāvṛṣi
when the earth has been excessively saturated with water, % atyarthañ jalopaklinnāyāṃ bhūmau
the accumulation [of wind], which is agitated by the cold, wind and rain, % śītavātavarṣerito 
produces wind diseases % vātikān vyādhīñ
in people whose bodies are extremely moist. % cātiklinnadehānāṃ dehināṃ

Thus, the causes of the accumulation and aggravation of these humours have been explained.

\item [12]
In this regard, the elimination of accumulated humours should occur in the monsoon, winter, and summer, and the elimination of aggravated humours in autumn, spring, and the dewy season. 
% note that the vulgate has nirharaṇaṃ kartavyam (kartavyam is missing in the Nepalese version), which indicates that eliminating these doṣas "should be done" in those seasons. 

\item [13]
Here, bile diseases are pacified in winter, phlegm diseases in summer, and wind diseases in autumn. 
Their accumulation, aggravation and pacification have been explained according to their nature. % svabhāvatas. 

\item [14]
Here, % the tatras are becoming tedious... should we start to leave them out?
in the forenoon of the day, % divasapūrvāhṇe : isn't divasa redundant here? The vulgate omits it.
one should observe % upalakṣayet
the characteristics of spring; % vasantaliṅgaṃ
at midday, those of summer; in the afternoon, those of the Prāvṛt monsoon; % see item 73 in 6.64 on the two rainy seasons (There, two rainy seasons called Prāvṛt and Varṣa are known to them)... 
in the evening, those of the Varṣa monsoon; 
at midnight, those of autumn; and at dawn, those of winter. %
Thus, one should know % evaṃ ... jānīyāt
the day and night like a year % ahorātram api varṣam iva
according to the cold, heat, rain, and the accumulation, aggravation and pacification of humours. % śītoṣṇavarṣadoṣopacayaprakopopaśamair

\item [15]
When the seasons are undisturbed, % avyāpanna : not disrupted, not unusual?
medicines are undisturbed, as is water. When they are consumed, they promote vitality, % prāṇa
longevity, % āyus 
strength, % bala
vigour, % vīrya
and \emph{ojas}. % I'm not sure how to translate ojas?

\item [16]
Harm to those medicines, however, is caused by unseen factors; % 'dṛṣṭakāritāni
the cold, heat, wind and rain.
As is known, % khalu
adverse [factors] % viparītāny ?
harm medicines and water as well. 

\item [17]
After taking them, % tāsām upayogād
various diseases manifest or an epidemic occurs. % maraka

\item [18]
Sometimes, % kadā cit
even when they are undisturbed, % avyāpanneṣu api
populations of people % janapadāḥ
are devastated % udhvasyante
through immoral actions caused by sorcery, % kṛtyā
curses, % abhiśāpa
demons, % rakṣaḥ
and anger; %krodhāt
or through the fragrance of poisonous herbs and flowers, % viṣauṣadhipuṣpagandhena
carried by the wind; % vāyunopanītena
or people are afflicted by cough, % kāsa
wheezing, % śvāsa
rheum, % pratiśyāya
head diseases, % śiroroga % Sharma has headache? I've seen śirovedanā for that though
and fever, 
or through the motion % carita
of planets and constellations; %grahanakṣatra
or through forbidden signs % garhitalakṣaṇa - Ḍalhaṇa seems to take this as forbidden signs (gṛhadāretyādi gṛhādīnāṃ yāni garhitāni ninditāni lakṣaṇāni cihnāni tannimittaṃ prādurbhāva utpattir yeṣāṃ vyādhīnāṃ tairvā upatapyante janapadāḥ) 
associated with beds, seats, conveyances, % yāna = dolārathādi (palanquins, chariots, etc.)
vehicles, % vāhanaṃ hastyaśvādi
jewels,
gems, % Ḍalhaṇa bangs on about the difference between maṇi and ratna. Doesn't seem important.
and vessels. % upakaraṇāni ghaṭapiṭharaśūrpādīni

\item [19]
When this occurs, % jāto 'tra
one should practise renunciation, % parityāga
rites of pacification and expiation, % śānti-prāyaścitta
offerings and auspicious observances, % bali-maṅgala
mantra recitation and fire rituals, % japa-homa
oblations and sacrifices, % upahāra-ijyā
\emph{añjali}s and salutations,% añjali-namaskāra
austerities and religious observances, % tapas-niyama
compassion and charity, % dayā-dāna
initiation and other such commitments to a teacher, % dīkṣā-abhyupagama: abhyupagama an agreement, contract? Ḍalhaṇa: abhyupagamaḥ guruvākyādīnām aṅgīkāraḥ; acceptance of the instructions of the guru, etc. note?
while remaining devoted to the deities, Brahmins, and the elders. % devatā-brāhmaṇa-guru-para
In this way, good arises. % evaṃ sādhu bhavati

\item [20]
And on this, the following verse:

The humours of people are aggravated during the seasons, whether the seasons have their normal qualities, have been disturbed, % viparīta
or are irregular. 
\end{translation} 